% =====================================================================
%  CHAPTER 5: 青春期发育（对应大纲第三章）
%  参考往届笔记：第五章 青春期生长发育
% =====================================================================
\chapter{青春期发育}
\setcounter{chapter}{5}
\chapterintro{
\textbf{【掌握】}青春期基本概念（青春期、青春发动期、青春发育进程、青春发动时相、\underline{青春发动时相长期变化}）；青春期内分泌变化（下丘脑、腺垂体、性腺等主要内分泌激素及HPG轴）；青春期体格发育特点（生长突增、身高速度高峰、两次交叉现象、身体各部分生长顺序、\underline{TOA/TOV/EA/EV四项核心指标}、三种成熟类型与五种生长模式）；\underline{性早熟（定义、中枢性/外周性分类、诊断标准、治疗原则）}；\underline{青春期发育迟缓（定义、CDGP、低促性腺激素型与高促性腺激素型、诊断流程与治疗原则）}。\par
\textbf{【熟悉】}青春期性发育的特殊表现（男女性器官、第二性征、性功能发育顺序及特点、\underline{Tanner性发育5阶段分期法}）；青春期性心理发育（\underline{性心理定义}、三个阶段、三大矛盾、\underline{性心理成熟四项标准}）。\par
青春期发动机制（中枢神经抑制机制减弱假说、HPG轴对性激素负反馈敏感性下降假说）。
}

% ----- 4.1 青春期与青春发动期 -----
\section{青春期与青春发动期}

\subsection{青春期基本概念}

\begin{keybox}
\textbf{一、青春期与青春发动期}

\textbf{1. 青春期（Adolescence）（WHO）：}个体从出现第二性征到性成熟的生理发育过程，是个体从儿童认知方式发展到成人认知方式的心理过程，是个体从经济的依赖性到相对独立状态的过渡，10-19岁。

\textbf{（1）青春发育进程（Pubertal Procession）：}描述某一时点，某个体在整个青春期发育过程中所处的绝对位置。特点：
\begin{enumerate}[leftmargin=*]
    \item 体格生长加速，以身高为代表的体格指标出现第二次生长突增。
    \item 各内脏器官体积增大、重量增加，功能日渐成熟。
    \item 内分泌功能活跃，与生长发育有关的激素分泌明显增加。
    \item 生殖系统功能发育骤然增快并迅速成熟，到青春晚期已具备繁殖后代的能力。
    \item 男女外生殖器和第二性征发育，使男女两性外部形态特征差别更明显。
    \item 在体格、功能发育同时，青春期心理发展加速，产生相应心理、行为变化，出现很多青春期特有的心理-行为问题。
\end{enumerate}

\textbf{（2）青春期分期：}早、中、晚三个阶段。
\begin{itemize}[leftmargin=*]
    \item \imp{青春早期：}主要表现为生长突增，出现身高突增高峰，性发育开始，一般持续约2-3年。
    \item \imp{青春中期：}以性器官、第二性征迅速发育为特征，出现月经初潮（女）或首次遗精（男），持续2-3年。
    \item \imp{青春后期：}体格生长速度明显减慢，但仍有增长，直至骨骺完全融合，性器官和第二性征持续发育至成人水平，社会心理发展加速，通常持续2年左右。
\end{itemize}

\vspace{3pt}
\begin{center}
\fbox{\parbox{0.95\textwidth}{
\textbf{【考点】}进入青春期第一信号：身高突增；性发育第一信号：男孩——睾丸增大，女孩——乳房发育；生殖功能开始成熟的标志：男孩——首次遗精，女孩——月经初潮。
}}
\end{center}
\vspace{3pt}

\textbf{2. 青春发动期（Puberty）：}青春期生殖系统发育与成熟的生物学变化过程，是青春期的早期阶段，时间跨度平均约4.5年（1.5-6年）。个体生理变化明显，包括体格生长突增、性腺和生殖器发育、第二性征出现等。

\textbf{（1）青春发动时相（Puberty Timing）：}在青春发动期，生殖系统发育的各种事件按特定时间模式进行，从生长突增（Growth Spurt）、乳房发育、腋毛生长、阴毛生长，到月经初潮等，这种青春发动期生殖系统发育的各种事件初现的时间模式。

\textbf{（2）}青春发动期各种事件初现时间存在明显个体差异，与群体平均水平比较可表现为提前（early）、适时（on-time）和推迟（delay）。一般以人群青春发动期生殖系统发育各种事件初现时间的前第25百分位数（P25）作为青春发动期时相提前的划界值。
\end{keybox}

\begin{defbox}
\textbf{3. 青春发动时相长期变化（Secular Change of Pubertal Timing）}

\imp{青春发动时相长期变化}是指青春发动期生殖系统发育各种事件出现时间在不同出生队列间呈现出的提前或延迟的长期趋势性变化现象。其中\imp{月经初潮年龄的长期下降}是最常使用的衡量指标。

\textbf{主要表现与机制：}
\begin{itemize}[leftmargin=*]
    \item 发达国家女童月经初潮年龄的下降速度已放缓或逐渐停止，而部分发展中国家仍呈下降趋势，揭示营养改善、社会经济水平提升对青春发动时相提前的推动作用。
    \item 青春发动时相长期变化受遗传背景、营养状况、体脂积累、社会环境、家庭经济水平以及环境内分泌干扰物（EDCs）等多因素的综合影响。儿童期超重肥胖是导致青春发动时相提前的重要危险因素。
    \item 青春发动时相的群体性提前与性早熟患病率上升密切相关，对儿童少年成年终身高、心理健康及远期生殖健康均产生深远影响。
\end{itemize}

\textbf{临床与公共卫生意义：}识别青春发动时相的长期变化趋势，有助于制定有针对性的青春期保健策略，对提前者加强监测与干预，对延迟者鉴别体质性与病理性原因。
\end{defbox}

% ----- 4.2 青春期内分泌变化 -----
\section{青春期内分泌变化}

\subsection{影响青春期发育的主要内分泌激素}

\begin{keybox}
\textbf{二、青春期内分泌变化}

\textbf{1. 影响青春期发育的主要内分泌激素}

\textbf{（1）下丘脑}
\begin{itemize}[leftmargin=*]
    \item 释放激素：促生长素释放激素、促甲状腺激素释放激素、促肾上腺皮质激素释放激素、促卵泡激素释放激素、黄体生成素释放激素、催乳素释放因子和促黑（素细胞）激素释放激素。
    \item 释放抑制激素：生长抑素、催乳素释放抑制因子和促黑（素细胞）激素释放抑制因子。
\end{itemize}

\textbf{（2）腺垂体}
\begin{itemize}[leftmargin=*]
    \item \imp{生长素（GH）}
    \begin{itemize}[leftmargin=*]
        \item 作用机制：a. 促使氨基酸进入细胞，加速蛋白质合成；b. 分解脂肪，使游离脂肪酸增加；c. 抑制葡萄糖氧化，减少糖原消耗。
        \item 调节：受下丘脑GHRH与GHRIH双重调节。GHRH促进生长素分泌，GHRIH抑制生长素分泌，同时外周生长素对下丘脑进行负反馈调节。生长素分泌与睡眠有关，入睡时分泌明显增加，入睡后60min左右血中生长素浓度达到高峰；慢波睡眠时相生长素分泌增加，快波睡眠时相分泌减少。
    \end{itemize}
    \item \imp{催乳素（PRL）}
    \begin{itemize}[leftmargin=*]
        \item 功能：促进乳腺发育，启动和维持乳腺泌乳。
        \item 调节：受下丘脑催乳素释放因子与催乳素释放抑制因子双重调节。正常生理情况下，后者对催乳素分泌的抑制作用占优势，所以青春期女性乳腺发育成熟、但无乳汁生成。
    \end{itemize}
    \item \imp{促激素}
    \begin{itemize}[leftmargin=*]
        \item 功能：调节各内分泌器官分泌相应内分泌激素。
        \item 分类：促甲状腺激素（TSH）、促肾上腺皮质激素（ACTH）、促卵泡激素（FSH）（促进卵泡成熟）和黄体生成素（LH）（促进排卵）。
    \end{itemize}
\end{itemize}

\textbf{（3）性腺：}产生精/卵子和分泌相应性激素。
\begin{itemize}[leftmargin=*]
    \item \imp{睾丸——雄激素：}睾酮（最重要）、双氢睾酮和脱氢表雄酮。睾酮的功能：
    \begin{itemize}[leftmargin=*]
        \item 胚胎发育期：对男性胎儿生殖器的分化起关键作用。
        \item 儿童期：促进体内蛋白质合成和骨骼、肌肉发育。
    \end{itemize}
    \item \imp{卵巢——雌激素：}\textbf{促进女性青春期乳腺发育最关键的激素。}
\end{itemize}

\textbf{（4）甲状腺：}甲状腺素（TH），功能：调节机体物质代谢、促进生长发育、维持神经和心血管功能。
\begin{itemize}[leftmargin=*]
    \item 出生后功能低下：呆小症/克汀病。
    \item 缺碘：地方性克汀病。
\end{itemize}

\textbf{（5）肾上腺：}肾上腺皮质与青春期生长发育有关。
\begin{itemize}[leftmargin=*]
    \item 雄激素：脱氢表雄酮，功能：调节性器官和第二性征发育。对女性阴毛、腋毛等第二性征的促进作用较明显，对男性性发育影响较小。
    \item 雌激素：肾上腺皮质分泌含量极少，对性发育影响不大。
\end{itemize}

\textbf{（6）胰岛：}胰岛$\beta$细胞分泌的胰岛素，功能：调节糖代谢，参与脂肪和蛋白质代谢（促进蛋白质合成和细胞增殖）。胰岛素与生长激素有协同作用，可与IGF-I受体结合而刺激生长。

\textbf{（7）松果体：}褪黑素，功能：抑制哺乳动物生殖系统发育。

\textbf{2. 青春期内分泌变化：}\imp{下丘脑-垂体-性腺（HPG）轴。}

\textbf{（1）性腺功能初现：}由下丘脑分泌的促性腺激素释放激素（GnRH）发动。男童：睾丸体积增大；女童：乳房发育。

\textbf{（2）肾上腺皮质功能初现：}由肾上腺分泌的肾上腺雄激素发动。出现时间较早（6-8岁），肾上腺雄激素分泌开始程序性增加，此时雄激素活性比睾酮作用弱，对男童作用更弱、显现不出作用；但对女童第二性征发育作用明显。
\end{keybox}

% ----- 4.3 青春期发动机制 -----
\section{青春期发动机制}

\subsection{神经-内分泌调控机制}

\begin{keybox}
\textbf{三、青春期发动机制}

\textbf{1. 神经-内分泌调控机制：}青春期内分泌调节的核心。

\textbf{2. 青春期发育的启动机制：}
\begin{enumerate}[leftmargin=*]
    \item 中枢神经抑制机制减弱；
    \item HPG轴对性激素负反馈敏感性下降。
\end{enumerate}
\end{keybox}

\subsection{启动机制详解}

\begin{keybox}
\textbf{（1）中枢神经抑制机制减弱（主流观点）：}

目前，多数学者认为，青春期的启动始于"青春期前静息期"这一下丘脑抑制机制的解除。从胎儿期开始，中枢神经对性发育就一直保持着抑制作用，使下丘脑对GnRH的分泌处于青春静止期（Juvenile Pause），内外生殖器官维持幼稚状态，而不继续发育。其后，伴随HPG轴的成熟，该抑制作用逐步减弱并消失，青春期开始启动。

青春期下丘脑GnRH脉冲释放是多因子参与的多层次网络性激发过程。由抑制性-兴奋性因子之间的平衡状态"调拨"启动性发育生物钟，这些因子包括神经递质、生长因子和某些神经肽类物质。在青春期前，抑制因子发挥主导作用抑制着GnRH的脉冲释放；当兴奋性因子功能增强或抑制因子作用减弱后，下丘脑GnRH脉冲释放增加，促使GnRH持续地释放，垂体对GnRH的敏感性提高，使腺垂体卵泡刺激素和黄体生成素等分泌量增加，从而启动青春发育。

\textbf{（2）HPG轴对性激素负反馈敏感性下降：}

童年期在外周血中已有少量可检测到的性激素，HPG系统对它们的存在极为敏感，受其负反馈抑制作用，HPG系统的活动处于静止状态。即将进入青春期时，下丘脑对这些外周激素的负反馈敏感性开始降低，分泌的GnRH逐渐增多，刺激垂体分泌FSH和LH，促进性腺发育和性激素分泌，青春期发育迅速推进。

\textbf{（3）Kisspeptin/GPR54系统——青春期启动的分子阀门：}

Kisspeptin是由Kiss-1基因编码的蛋白质，G蛋白偶联受体54（GPR54）是其特异性受体。GPR54/Kisspeptins在中枢神经系统尤其是下丘脑有丰富的表达，\imp{是青春期发育启动的分子阀门}，是激活GnRH神经元的一个重要的兴奋性因素，其调控机制复杂，受PcG、甲基化表观遗传等多种因素调节。如makorin环指蛋白3（MKRN3）基因在童年期活性高，选择性地抑制Kiss-1和TAC3的启动子活性，减少kisspeptin和NKB的释放并减少GnRH分泌；而到了青春发动期，受miR-30家族成员的调节，MKRN3基因活性下降，对GnRH和NKB及Kiss-1基因有抑制解除，GnRH分泌增加，启动了性腺功能发育。
\end{keybox}

% ----- 4.4 青春期生长发育的特点 -----
\section{青春期生长发育的特点}

\subsection{青春期体格发育特点}

\begin{keybox}
\textbf{一、青春期体格发育特点}

\textbf{1. 青春期生长突增}

\textbf{（1）生长突增（Growth Spurt）：}进入青春期后最引人注目的形态变化是儿童少年体格生长出现的突发性快速生长现象。

\begin{itemize}[leftmargin=*]
    \item \imp{身高速度高峰（Peak Height Velocity, PHV）：}以身高为代表，生长速度在童年期较平稳的每年4-5cm基础上出现快速增长，1-2年后达到高峰。PHV的发生时的年龄称\imp{身高速度高峰年龄（Age at Peak Height, PHA）}。
    \item \imp{体重速度高峰（Peak Weight Velocity, PWV）和体重速度高峰年龄（Age at Peak Weight, PWA）：}儿童少年体格发育的另一重要指标，在青春期发育阶段的生长发育特点与身高发育特点基本一致。
\end{itemize}

\textbf{2. 青春期生长突增的速度变化}

\textbf{（1）}年龄别生长速度曲线和按PHA生长速度曲线，存在明显不同。
\begin{itemize}[leftmargin=*]
    \item 由于青春期体格生长突增高峰年龄的个体变异突出，按年龄别生长速度曲线中儿童身高生长突增高峰值明显小于按PHA生长速度曲线中儿童身高生长突增高峰值；按年龄别生长速度曲线波动性相对更小，曲线更像是经平滑统计处理后的生长速度曲线。
    \item 因此，评价个体儿童青春期生长发育速度时，不能单独应用群体按年龄生长速度参考值，可将按年龄生长速度参考值和按PHA生长速度参考值结合考虑。
\end{itemize}

\textbf{（2）男女身高、体重生长曲线出现两次交叉现象。}
\begin{itemize}[leftmargin=*]
    \item \imp{第一次交叉：}女性因体格生长突增开始早，青春期开始早期女性体格生长水平超过同龄男性。
    \item \imp{第二次交叉：}随年龄增长女性体格生长逐步进入缓慢阶段，而男性体格生长突增开始，在男性体格生长速度达到高峰前，其生长水平已超过女性。
\end{itemize}

\textbf{（3）青春期身体各部分的生长速度：}由于青春期身体各部分生长突增不是同时开始的，生长速度也不相同，故青春期间身体各部分的比例不断变化。
\begin{itemize}[leftmargin=*]
    \item 青春前期，上下肢增长早于躯干，下肢增长稍早于上肢。上下肢生长中，顺序大致为：足长→下肢长→手长→上肢长；下肢生长中，足长生长首先加速，也最早停止；一般14-15岁后足长即不再增长。足长加速生长后6个月，小腿开始增长，然后是大腿。由于下肢增长早，青春前期坐高指数（坐高/身高）逐渐下降，中期降至最低点；7岁时为55.2（男）和55.0（女），13岁时降至52.9（男）53.4（女），因而出现青春早期长臂长腿体态。当小腿增长达顶点后4个月，骨盆宽、胸宽开始增长，11个月后肩宽增长加速。青春中后期，躯干增长速度加快，故坐高指数再度增加，成年时约54.2（男）和54.3（女）。
\end{itemize}

\vspace{3pt}
\begin{center}
\fbox{\parbox{0.95\textwidth}{
\textbf{【考点】}可根据青少年足长最先突增又较早停止生长的特点，用足长预测成年身高。
}}
\end{center}
\vspace{3pt}

\begin{itemize}[leftmargin=*]
    \item 青春期体格生长突增不仅体现在身高、体重等形态指标方面，也体现在体内组织器官上；青春期生长突增无论在男女间，在身体各部分间，都存在突增时间和幅度差异。
\end{itemize}

\textbf{3. 青春期生长发育类型}

\textbf{（1）3种成熟类型生长特点}
\begin{itemize}[leftmargin=*]
    \item \imp{早熟型：}发育过程中体重/身高比值一直大于晚熟型，骨盆较宽、肩部较窄，最后形成盆宽、窄肩的相对矮胖体型。多为高度女性体型特征，早熟型男孩体型相对偏向女性。青春期启动早，女孩8-9岁，男孩10-11岁。生长突增高峰的出现和停止时间都早，生长突增时间持续一年左右，生长期较短。女孩早熟型相对多于男孩。
    \item \imp{晚熟型：}发育过程中体重/身高比值一直小于早熟型，骨盆较窄、肩部较宽，最后形成盆窄、肩宽的相对瘦高体型。多为高度男性体型特征，晚熟型女孩体型相对偏向男性。青春期启动晚，女孩10-11岁、男孩12-13岁。生长突增高峰的出现和停止时间都晚，突增维持时间较长，可达2年以上甚至3年，生长期较长。男孩晚熟型相对多于女孩。
    \item \imp{一般型：}发育过程中体重/身高比值、肩宽和盆宽、青春期启动年龄和体型特征等，都介于早熟型和晚熟型间。青春期生长突增时间多维持在2年左右。
\end{itemize}

\textbf{（2）成熟类型与生长突增——五种生长模式：}
\begin{enumerate}[leftmargin=*]
    \item 模式I：生长突增起始于平均年龄，成年身高位于平均水平。
    \item 模式II：生长突增发生较早，突增时身高大于同龄者，突增结束年龄也较早，导致其生长期较短、身高增长量较少，成年身高低于平均水平。
    \item 模式III：儿童期和青春早期生长都低于同龄者，但较晚的生长突增和较长的生长期导致其成年身高达到甚至超过平均水平。
    \item 模式IV：遗传性（家族性）高身材。
    \item 模式V：遗传性（家族性）矮身材。
\end{enumerate}

\vspace{3pt}
\begin{center}
\fbox{\parbox{0.95\textwidth}{
\textbf{【考点】}后两种生长模式（模式IV和模式V）尽管其生长突增起始于平均年龄，生长速度也较相似，但因遗传影响，其平均身高始终分别处于高水平或低水平。
}}
\end{center}
\vspace{3pt}
\end{keybox}

\subsection{青春期性发育特点}

\begin{defbox}
\textbf{二、青春期性发育特点}

\textbf{定义：}青春期性发育是青春期发育最重要的特征之一，包括内外生殖器官的变化、第二性征的发育、生殖功能的发育成熟等。
\end{defbox}

\begin{keybox}
\textbf{1. 男性性发育}

\textbf{（1）性器官形态发育：}男孩青春期性发育存在很大个体差异，但各指征的出现顺序大致相似：睾丸→阴茎、身高突增（一年后）。

\textbf{（2）第二性征发育：}阴毛→腋毛→胡须→变声、喉结出现。阴毛一般11-12岁左右出现，1-2年后出现腋毛，再隔1年左右胡须开始萌出，额部发际后移，脸型轮廓从童年型向成年型演变。随着体内雄激素水平增高，喉结增大，声带变厚变长，一般13岁左右出现变声。绝大多数男孩在18岁前完成第二性征发育。

\vspace{3pt}
\begin{center}
\fbox{\parbox{0.95\textwidth}{
\textbf{【考点】}第二性征（Secondary Sexual Characteristics）：人或动物发育后表现出的与生殖系统无直接关系、而与性别有关的机体外表特征。
}}
\end{center}
\vspace{3pt}

\textbf{（3）性功能发育：}随着睾丸生长，生殖功能开始逐渐发育成熟。遗精是男性生殖功能开始发育成熟的重要标志之一。首次遗精一般发生于12-18岁间。随后身高发育速度逐步减慢，而睾丸、附睾、阴茎等迅速发育，逐步接近成人水平。

\textbf{2. 女性性发育}

\textbf{（1）性器官形态发育：}进入青春期后，在垂体/性腺分泌的促卵泡激素、黄体生成素、促性腺激素作用下，内外生殖器迅速发育。
\begin{itemize}[leftmargin=*]
    \item 卵巢：增重，发育成熟；
    \item 子宫：增重增长，宫体增长大于宫颈；
    \item 外生殖器：阴阜因脂肪堆积隆起，大阴唇变厚、小阴唇变大、有色素沉着，出现阴道分泌物、碱性→酸性。
\end{itemize}

\textbf{（2）第二性征发育：}乳房→阴毛→腋毛发育。乳房发育平均开始于11岁，从乳房发育I-V度，历时约4年。乳房开始发育后6个月-1年出现阴毛；而后6个月-1年出现腋毛。身高生长突增几乎与乳房发育同时或稍早开始，而身高突增高峰出现年龄通常在乳房发育后1年左右。

\textbf{（3）性功能发育：}月经初潮（在身高突增高峰后1年左右）。
\end{keybox}

\subsection{Tanner性发育5阶段分期法}

\begin{keybox}
\textbf{Tanner青春期性发育5阶段分期法（Tanner Stages / Sexual Maturity Rating, SMR）}

Tanner分期是由英国儿科医生James Tanner提出的国际通用青春期性发育评估标准，通过对\imp{男童外生殖器}和\imp{女童乳房}的形态变化进行五级分期，客观量化性发育进程，是青春期生长发育评价和临床诊疗中最核心的性发育评估工具。

\textbf{男童外生殖器Tanner分期（按睾丸容积、阴囊和阴茎形态）：}
\begin{enumerate}[leftmargin=*]
    \item \imp{Tanner I期（青春前期）：}睾丸容积$<$4 ml，阴茎和阴囊仍维持儿童早期的大小与比例，呈幼稚型。
    \item \imp{Tanner II期：}睾丸容积增大至4--8 ml，阴囊皮肤颜色变红、纹理改变；阴茎无变化或变化很小。此期标志\imp{性腺功能初现（Gonadarche）}启动。
    \item \imp{Tanner III期：}睾丸容积进一步增大至8--12 ml，阴茎长度增加，睾丸和阴囊继续增大。
    \item \imp{Tanner IV期：}睾丸容积12--15 ml，阴茎头增粗发育，阴茎进一步增大，龟头显露；阴囊皮肤颜色加深。
    \item \imp{Tanner V期（成人期）：}外生殖器大小和形状达到成人水平，睾丸容积$\geq$15 ml。
\end{enumerate}

\textbf{女童乳房Tanner分期：}
\begin{enumerate}[leftmargin=*]
    \item \imp{Tanner I期（青春前期）：}发育前期，仅有乳头突出，无乳腺组织可触及。
    \item \imp{Tanner II期（乳腺萌出期）：}乳腺隆起，乳房和乳晕呈单个小丘状隆起，乳晕增大。此期标志\imp{性腺功能初现（Gonadarche）}启动，通常是女童进入青春期的第一个可见信号。
    \item \imp{Tanner III期：}乳房和乳晕进一步增大，但两者仍在同一个丘状水平面上，乳晕色素加深。
    \item \imp{Tanner IV期：}乳头和乳晕突出于乳房丘面上，形成第二个小丘，为女童性发育的特有阶段。
    \item \imp{Tanner V期（成熟期）：}乳房增大至成人形态，乳晕回落到乳房整体轮廓之内，乳房和乳晕又在同一个丘面上。
\end{enumerate}

\vspace{3pt}
\begin{center}
\fbox{\parbox{0.95\textwidth}{
\textbf{【考点】}II期为性腺功能初现标志——男童睾丸$\geq$4 ml，女童乳房萌出。女童IV期（乳晕形成第二小丘）是女性青春期特有标志，男童无对应阶段。
}}
\end{center}
\vspace{3pt}

\textbf{临床应用要点：}
\begin{itemize}[leftmargin=*]
    \item Tanner分期是评价性发育进程的\imp{金标准}，也是性早熟和青春期发育迟缓诊断中判断性发育水平的核心依据。
    \item 阴毛发育（Pubarche）虽常与Tanner分期并行评估，但主要受\imp{肾上腺皮质功能初现（Adrenarche）}驱动，可与性腺功能初现分离，故不纳入Tanner分期核心指标。
    \item 分期评估应通过视诊和触诊进行，避免仅凭受检者自评。
\end{itemize}
\end{keybox}

\subsection{青春期性心理发育}

\begin{defbox}
\textbf{三、青春期性心理发育}

\textbf{性心理定义：}\imp{性心理（Sexual Psychology）}是指围绕性特征、性欲望和性行为而展开的所有心理活动，是由性意识、性感情、性知识、性经验、性观念等要素结构而成的多层次心理复合体。性心理是个体人格发展的重要组成部分，在青春期经历从萌芽、发展到成熟的关键过渡。

\textbf{1. 发展阶段}
\begin{enumerate}[leftmargin=*]
    \item \imp{异性疏远期：}性发育早期体态剧烈变化使青少年心理出现许多不同心理现象，包括心理适应困难，内心慌乱不安，表现出对性的抵触现象，甚至反感。大部分男女生这时开始相互疏远异性，表现出男女排斥现象。
    \item \imp{异性爱慕期：}随着青春期性心理逐步发展，男女都开始渴望了解异性，主动接近异性，同时开始萌发对异性的向往和憧憬。
    \item \imp{两性初恋期：}性意识的发展逐渐趋于成熟，对异性的情感充满浪漫和幻想，对异性爱慕逐渐趋于专一。在两性初恋期，男女生情感都不稳定，并有明显占有心理。
\end{enumerate}

\textbf{2. 青春期性心理发展矛盾现象}
\begin{enumerate}[leftmargin=*]
    \item 性生理发育成熟与性心理相对幼稚的矛盾。
    \item 自我意识迅猛发展与社会成熟度相对迟缓的矛盾。
    \item 情感激荡释放与外部表露趋向内隐的矛盾。
\end{enumerate}

\textbf{3. 性心理成熟的一般标准}
\begin{enumerate}[leftmargin=*]
    \item \imp{认知标准：}能够正确理解男女两性差异和性关系的内涵，正确领悟性的生物学意义和社会意义。
    \item \imp{情感标准：}在心理上具有适度的性需要，内心乐意与异性进行正常接触和交往，能够建立健康的爱情观和恋爱关系。
    \item \imp{意志标准：}具有正常的性意识和较强的性行为理智感，能够自觉运用社会规范和自我意志控制性行为冲动。
    \item \imp{社会适应标准：}在心理上具备适应正常婚姻状态的能力，能够承担与性相关的社会角色和责任。
\end{enumerate}

\vspace{3pt}
\begin{center}
\fbox{\parbox{0.95\textwidth}{
\textbf{【考点】}青春期是性心理障碍开始表现的关键期。
}}
\end{center}
\vspace{3pt}
\end{defbox}

\subsection{性早熟}

\begin{defbox}
\textbf{一、性早熟（Precocious Puberty）}

\imp{性早熟}是指女童在7.5岁前出现乳房发育或10.0岁前出现月经初潮，男童在9.0岁前出现睾丸发育（睾丸容积$\geq$4 ml，达Tanner II期）。其核心危害在于性激素过早暴露导致骨骺提前闭合、损害成年终身高，以及性征早现带来的心理行为问题和社会适应困难。

\textbf{分类（按发病机制）：}
\begin{enumerate}[leftmargin=*]
    \item \imp{中枢性性早熟（CPP）：}又称GnRH依赖性性早熟或真性性早熟，是由于HPG轴功能提前启动，下丘脑脉冲性释放GnRH增加，导致垂体促性腺激素分泌增多，继而刺激性腺发育并分泌性激素，使内外生殖器发育和第二性征按正常程序呈现。CPP的性发育过程遵循正常的青春期发育程序，具有生育能力。
    \item \imp{外周性性早熟（PPP）：}又称非GnRH依赖性性早熟或假性性早熟，是由各种原因（性腺/肾上腺肿瘤或增生、外源性性激素摄入等）引起的体内性甾体激素升高至青春期水平，导致第二性征早现，但不依赖GnRH的驱动。PPP仅有第二性征的部分早现，\imp{不具有完整的性发育程序性过程}，性腺维持在青春前期状态，不具备生殖能力。
\end{enumerate}
\end{defbox}

\begin{keybox}
\textbf{1. 中枢性性早熟（CPP）}

\textbf{病因：}
\begin{itemize}[leftmargin=*]
    \item \imp{特发性CPP：}未能发现明确的器质性病变，约占女童CPP的80\%--90\%，男童CPP的40\%--50\%。
    \item \imp{中枢神经系统器质性病变：}下丘脑错构瘤（最常见的CPP器质性病因）、其他下丘脑或垂体肿瘤、脑外伤、脑炎、脑积水、蛛网膜囊肿、放射治疗后等。
    \item \imp{继发性CPP：}由外周性性早熟长期未控制，外周性激素水平持续升高，经负反馈机制激活HPG轴转化而来。
\end{itemize}

\textbf{CPP诊断标准（2022年中国专家共识）：}
\begin{enumerate}[leftmargin=*]
    \item 性征提前出现：女童7.5岁前乳房发育，男童9.0岁前睾丸增大$\geq$4 ml。
    \item 性腺增大：女童盆腔B超见卵巢容积$>$1 ml，并可见多个直径$>$4 mm的卵泡；男童睾丸容积$\geq$4 ml。
    \item 血清促性腺激素及性激素达青春期水平：GnRH激发试验是诊断CPP的\imp{金标准}——以GnRH 2.5 $\mu$g/kg（最大100 $\mu$g）静脉推注后，免疫化学发光法检测LH峰值$\geq$5.0 IU/L且LH峰/FSH峰$\geq$0.6，提示HPG轴已激活。
    \item 多有骨龄提前，骨龄超过实际年龄$\geq$1岁。
    \item 有线性生长加速，年生长速度高于同龄健康儿童。
\end{enumerate}

\textbf{治疗原则：}GnRH类似物（GnRHa）是当前CPP治疗的一线选择，通过持续作用于垂体GnRH受体使受体下调而抑制HPG轴功能，使性发育暂停或逆转，核心目标是改善因骨龄提前而减损的成年终身高。器质性病因所致CPP需同时针对原发病进行治疗。

\vspace{6pt}

\textbf{2. 外周性性早熟（PPP）}

\textbf{病因：}
\begin{itemize}[leftmargin=*]
    \item \imp{同性性早熟（与原性别一致）：}女童——McCune-Albright综合征、自律性卵巢囊肿、分泌雌激素的肿瘤、外源性雌激素摄入；男童——先天性肾上腺皮质增生症（CAH，最常见）、肾上腺皮质肿瘤或睾丸间质细胞瘤、外源性雄激素摄入。
    \item \imp{异性性早熟（与原性别相反）：}女童——CAH（21-羟化酶缺乏最常见）、分泌雄激素的肿瘤、外源性雄激素摄入；男童——分泌雌激素的肿瘤、外源性雌激素摄入。
\end{itemize}

\textbf{PPP诊断要点：}第二性征提前出现但性征发育不按正常青春发育程序进展；性腺大小维持在青春前期水平；促性腺激素在青春前期水平，GnRH激发试验阴性（LH峰$<$5.0 IU/L，LH峰/FSH峰$<$0.6）。

\textbf{治疗原则：}按不同病因分别处理——各类肿瘤行手术切除；CAH予以糖皮质激素替代治疗，抑制ACTH过度分泌；外源性激素摄入者终止暴露源。
\end{keybox}

\subsection{青春期发育迟缓}

\begin{defbox}
\textbf{二、青春期发育迟缓（Delayed Puberty）}

\imp{青春期发育迟缓}是指在正常青春期启动平均年龄2.5个标准差以上（女童$\geq$13岁，男童$\geq$14岁）仍无第二性征发育的征兆，或从青春期启动至发育完成的时间超过5年。在普通青少年人群中的发生率约2\%--3\%，男童多于女童。主要分为\imp{体质性青春期生长发育延迟（CDGP）}和\imp{青春期发育障碍（病理性）}两大类。
\end{defbox}

\begin{keybox}
\textbf{1. 体质性青春期生长发育延迟（CDGP）}

CDGP是青春期发育迟缓中最常见的类型，约占全部病例的2/3。其本质是一种\imp{边缘性生理现象}，属于正常生长发育的极端变异，而非器质性疾病。

\textbf{发病机制：}HPG轴调控机制的成熟滞后——下丘脑GnRH脉冲发生器的激活时间延迟，导致促性腺激素和性激素水平的升高晚于同龄人，但最终的性发育和生殖功能完全正常。常有\imp{家族遗传倾向}（父母一方或双方有"晚发育"史）。

\textbf{临床特征：}出生体重和身长正常，儿童期生长速度处于正常低限，骨龄落后于实际年龄（但骨龄与身高年龄相符）；青春期启动延迟，但一旦启动后性发育过程正常，最终可实现完全正常的性成熟和成年身高。属于\imp{暂时性自限性状态}。

\textbf{处理原则：}以解释、安抚和定期随访监测为主，配合心理支持措施。少数心理负担极重者可短期小剂量性激素诱导青春期启动。

\vspace{6pt}

\textbf{2. 青春期发育障碍（病理性）}

约占青春期发育迟缓病例的1/3，根据促性腺激素水平分为两型：

\textbf{（1）低促性腺激素型性腺功能减退（HH）——继发性/中枢性：}
\begin{itemize}[leftmargin=*]
    \item \imp{发病机制：}因先天性或后天性原因导致下丘脑和/或垂体病变，引起GnRH或LH/FSH的生成和分泌减少，继而导致性腺因缺乏上游激素刺激而不能正常发育。
    \item \imp{激素特征：}血清LH、FSH水平降低或处于正常低值，性激素水平降低。
    \item \imp{主要病因：}先天性——Kallmann综合征（GnRH神经元迁移障碍伴嗅觉缺失）、特发性IHH、多种垂体激素缺乏症等；后天性——下丘脑或垂体肿瘤（颅咽管瘤最常见）、颅脑外伤或手术、放射治疗后、功能性促性腺激素缺乏（神经性厌食、过度运动、慢性消耗性疾病）。
\end{itemize}

\textbf{（2）高促性腺激素型性腺功能减退——原发性/外周性：}
\begin{itemize}[leftmargin=*]
    \item \imp{发病机制：}因先天性或后天性原因导致睾丸或卵巢本身的病变，引起性激素生成和分泌减少，通过负反馈机制使垂体LH及FSH\imp{继发性显著升高}。
    \item \imp{激素特征：}血清LH、FSH水平明显升高（FSH升高尤为突出），性激素水平降低。
    \item \imp{主要病因：}先天性——Turner综合征（45,X——女童最常见）、Klinefelter综合征（47,XXY——男童最常见）、性腺发育不良等；后天性——睾丸或卵巢外伤/扭转/手术切除、感染（腮腺炎性睾丸炎）、放化疗后性腺损伤、自身免疫性卵巢炎。
\end{itemize}

\vspace{6pt}

\textbf{3. 诊断流程：}核心任务是区分CDGP与病理性发育障碍。
\begin{enumerate}[leftmargin=*]
    \item 详细病史：出生史、生长发育史、青春期发育家族史（父母"晚发育"史提示CDGP）、嗅觉功能（筛查Kallmann综合征）。
    \item 全面体格检查：生长曲线评估、Tanner分期评定、有无Turner/Klinefelter综合征的特殊体征。
    \item 骨龄测定：骨龄显著落后于实际年龄是CDGP的重要特征。
    \item 实验室检查：第一线——血清LH、FSH、睾酮（男）/雌二醇（女）、甲状腺功能、IGF-I；第二线——GnRH激发试验、hCG激发试验（男童）；第三线——染色体核型分析、头颅MRI。
\end{enumerate}

\textbf{4. 治疗原则：}
\begin{itemize}[leftmargin=*]
    \item \imp{心理支持与辅导：}所有类型患儿的共同基础治疗，减轻因发育延迟带来的焦虑、自卑和社交回避。
    \item \imp{病因治疗：}确诊的肿瘤行手术切除或综合治疗；功能性促性腺激素缺乏者纠正原发因素。
    \item \imp{激素替代治疗：}确诊的病理性性腺功能减退采用性激素替代疗法（男童：低剂量睾酮递增方案；女童：低剂量雌激素递增方案，后期联合孕激素建立人工月经周期）。
\end{itemize}

\vspace{3pt}
\begin{center}
\fbox{\parbox{0.95\textwidth}{
\textbf{【考点】}CDGP与病理性HH的鉴别要点：CDGP有"晚发育"家族史、骨龄与身高年龄一致、GnRH激发试验呈青春前期反应模式；HH常无家族晚发育史、GnRH激发试验呈低平曲线。高促性腺激素型以LH/FSH显著升高为特征，易于与前两者鉴别。
}}
\end{center}
\vspace{3pt}
\end{keybox}

\newpage
